% !TeX program = xelatex
\documentclass[12pt,a4paper]{exam-zh}

\usepackage{amsmath,array,graphicx,booktabs,multirow}
\usepackage{enumitem}
\usepackage{fancyhdr}
\usepackage{tikz}
\usepackage{tabularx}
\usepackage{setspace}
\setCJKmainfont{Noto Serif CJK SC}
\geometry{margin=2.2cm,top=2.0cm,bottom=2.2cm}
\pagestyle{fancy}
\fancyhf{}
\cfoot{-\thepage{}-}
\renewcommand{\headrulewidth}{0pt}
\renewcommand{\footrulewidth}{0pt}
\setlength{\parindent}{0pt}
\setlength{\parskip}{0.2em}
\setlength{\headheight}{14pt}
\linespread{1.12}
\usepackage{setspace}
\usepackage{eso-pic}
\usetikzlibrary{calc}

%region exam setup

\examsetup{
    page = {
    size = a4paper,           % 设定纸张大小 [4]
    foot-content = { - ; - }  % 页脚显示为 -1-, -2- 格式 [3]
  },
    choices = {
    label-pos = top-left,
    column-sep = 1em,
    label-pos = auto,
    label-sep = 0.5em,
    max-columns = 4
    },
  question = {
    top-sep    = 1em, % 调整题目上方的垂直距离（默认值为 .5em plus .5em minus .2em）
    bottom-sep = 1em  % 调整题目下方的垂直距离（默认值为 .5em plus .5em minus .2em）
  },
    paren = {
    show-paren = false  % 设置为 false 则不显示括号 [3, 4]
  },
    choices = {
    label = \textbf{\Alph*}
  }, 
    solution = {
    show-solution = show-move, % 答案自动移动到卷末
    label-content = {参考答案},  % 答案区的标题文字
    show-qed = true            % 在解答末尾显示结束符号 ◻ 
  }
}

%endregion


%region 卷面变量设置
    % ==========================
    % 兼容从当前文件夹或工作区根目录进行编译
    % ==========================
% 卷面变量设置（只需改这里）
% ==========================
\newcommand{\papercode}{SC 07}
\newcommand{\examtitle}{2026年金宝培元独立中学毕业考}
\newcommand{\examlevel}{高中组}
\newcommand{\examsubject}{高级数学(I)}
\newcommand{\sidebartext}{高级数学(I)/一}
\newcommand{\papername}{试卷一}
\newcommand{\papertype}{选择题}
\newcommand{\examdate}{2026年9月24日}
\newcommand{\examtime}{8.00 a.m.－9.00 a.m.}
\newcommand{\examduration}{60分钟}
\newcommand{\questionpages}{4面}
\newcommand{\paperonepercentage}{30\%}
\newcommand{\papertwopercentage}{70\%}
\newcommand{\setterteacher}{蔡杰庆师}
\newcommand{\reviewteacher}{杨丽芳师}
\newcommand{\auditteacher}{雷泳健师}


    \newcommand{\optrow}[4]{%
    \begin{tabular*}{\linewidth}{@{\extracolsep{\fill}}llll}
    A~#1 & B~#2 & C~#3 & D~#4
    \end{tabular*}%
    }
    \newcommand{\qitem}[5]{\item #1\par\optrow{#2}{#3}{#4}{#5}}
%endregion
    \begin{document}


%region 封面
    \AddToShipoutPictureFG*{%
    \begin{tikzpicture}[remember picture,overlay]
        \node[
            anchor=north west,
            fill=black,
            minimum width=22pt,
            minimum height=90pt,
            inner sep=0pt
        ] (sidebar) at ([xshift=3em,yshift=-3em]current page.north west) {};
        \node[
            rotate=270,
            text=white,
            font=\bfseries\footnotesize
        ] at (sidebar.center) {\sidebartext};
    \end{tikzpicture}%
    \put(24,18){\makebox(0,0)[lb]{\small 版权所有}}%
    \put(520,18){\makebox(0,0)[lb]{\bfseries\small 【机密】}}%
    }

    \thispagestyle{fancy}

    \begin{minipage}[t]{0.22\linewidth}
    \vspace{0pt}
    \end{minipage}
    \hfill
    \begin{minipage}[t]{0.8\linewidth}
        \centering
        {\zihao{2}\bfseries \examtitle\par}
        \vspace{1.3em}
        {\zihao{3}\bfseries \examlevel\par}
        \vspace{1.3em}
        {\zihao{2}\bfseries \examsubject\par}
        \vspace{0.2em}
        {\large (\papercode)\par}
    \end{minipage}
    \hfill
    \begin{minipage}[t]{0.18\linewidth}
    \vspace{0pt}
    \end{minipage}

    \vspace{3.0em}

    \begin{center}
    \begin{tabular}{l}
        \multicolumn{1}{c}{\papername\hspace{2em}\papertype}\\[-0.1em]
        \midrule
        日期：\examdate\\[-0.1em]
        时间：\examtime\\[-0.1em]
        \multicolumn{1}{c}{(\examduration)}\\
        \midrule
    \end{tabular}
    \end{center}

    \vspace{2.0em}

    {\zihao{4}\bfseries \underline{考生须知}\par}
    \vspace{0.8em}

    \begin{enumerate}[label=（\chinese*）,leftmargin=3em,itemsep=0.65em]
    \item 本科试卷共分两份：\\
    \hspace*{2em}试卷一：选择题（\paperonepercentage）；\\
    \hspace*{2em}试卷二：作答题（\papertwopercentage）。

    \item 考生须于第一阶段规定的60分钟内完成试卷一。第二阶段规定的2小时内作答试卷二。

    \item 试卷一选择题15题全答。选出正确的答案，然后用\textbf{5B铅笔将“O”答案纸}（电脑卡）上相应的字母所在的小方格涂黑。

    \item 可使用电子计算机进行演算。

    \item 试卷一选择题共印\questionpages。
    \end{enumerate}

    \vspace{2.0em}

    \begin{center}
    {% 外框粗细及外、内框之间的距离
    \setlength{\fboxrule}{3pt}%
    \setlength{\fboxsep}{1.5pt}%
    \fbox{%
        {% 内框粗细及文字与内框之间的距离
        \setlength{\fboxrule}{0.5pt}%
        \setlength{\fboxsep}{6pt}%
        \fbox{\rule{0pt}{1.1em}\hspace{1.2em}{\zihao{4}\bfseries 未经正式宣布\quad 不得翻看内页}\hspace{1.2em}}%
        }%
    }%
    }
    \end{center}

    \vfill

    \begin{tabular*}{\linewidth}{@{\extracolsep{\fill}}ll}
    考生姓名：\underline{\hspace{4.2cm}} & 考生编号：\underline{\hspace{4.0cm}}
    \end{tabular*}

%endregion

\newpage





%region 范围：2026年度金宝培元独中毕业考·高中组·高级数学试卷一 | 考点：多项式整除与待定系数法 | 难度：基础
\begin{question}
    已知多项式 $x^{2}-x$ 整除 $x^{4}-3x^{3}+3x^{2}-mx$，求 $m$ 的值。\paren[A]
    \begin{choices}
        \item $1$
        \item $2$
        \item $3$
        \item $4$
    \end{choices}
\end{question}

\begin{solution}
    已知 $x^{2}-x$ 整除 $x^{4}-3x^{3}+3x^{2}-mx$，求 $m$。 \\
    设商式为 $x^{2}+ax+b$，用待定系数法比较系数。
    \begin{align*}
        x^{4}-3x^{3}+3x^{2}-mx &= (x^{2}-x)(x^{2}+ax+b) \\
        &= x^{4}+(a-1)x^{3}+(b-a)x^{2}-bx
    \end{align*}
    比较系数：
    \begin{align*}
        a-1 &= -3 \ \Rightarrow \ a=-2 \\
        b-a &= 3 \ \Rightarrow \ b=1 \\
        m &= b = 1
    \end{align*}
    故选 A。
\end{solution}
%endregion

% %region 范围：2026年度金宝培元独中毕业考·高中组·高级数学试卷一 | 考点：一元二次方程根与系数关系（韦达定理） | 难度：基础
% \begin{question}
%     设 $\alpha$ 及 $\beta$ 为方程 $2x^{2}+kx+1=0$ 的两个根，试用 $\alpha,\beta$ 表示 $k$。\paren[A]
%     \begin{choices}
%         \item $-2(\alpha+\beta)$
%         \item $-\alpha-\beta$
%         \item $2(\alpha+\beta)$
%         \item $\alpha+\beta$
%     \end{choices}
% \end{question}

% \begin{solution}
%     已知 $\alpha,\beta$ 为 $2x^{2}+kx+1=0$ 的两根，求 $k$。 \\
%     由韦达定理，两根之和等于 $-\dfrac{k}{2}$。
%     \begin{align*}
%         \alpha+\beta &= -\frac{k}{2} \\
%         k &= -2(\alpha+\beta)
%     \end{align*}
%     故选 A。
% \end{solution}
% %endregion

%region 范围：2026年度金宝培元独中毕业考·高中组·高级数学试卷一 | 考点：等差数列通项公式与求和公式 | 难度：中等
\begin{question}
    一等差数列的首四项之和为 $\dfrac{17}{3}$，第十项为 $\dfrac{31}{6}$，求该数列的公差。\paren[D]
    \begin{choices}
        \item $\dfrac{7}{6}$
        \item $\dfrac{5}{3}$
        \item $\dfrac{2}{3}$
        \item $\dfrac{1}{2}$
    \end{choices}
\end{question}

\begin{solution}
    已知首四项和为 $\dfrac{17}{3}$，第十项为 $\dfrac{31}{6}$，求公差 $d$。 \\
    设首项为 $a_{1}$，公差为 $d$，分别列出两个条件的方程联立求解。
    \begin{align*}
        4a_{1}+6d &= \frac{17}{3} \\
        a_{1}+9d &= \frac{31}{6}
    \end{align*}
    由第一式：
    \[
        a_{1} = \frac{17}{12}-\frac{3d}{2}
    \]
    代入第二式：
    \begin{align*}
        \frac{17}{12}-\frac{3d}{2}+9d &= \frac{31}{6} \\
        \frac{15d}{2} &= \frac{31}{6}-\frac{17}{12} = \frac{15}{4} \\
        d &= \frac{1}{2}
    \end{align*}
    故选 D。
\end{solution}
%endregion

%region 范围：2026年度金宝培元独中毕业考·高中组·高级数学试卷一 | 考点：数列求和公式（平方和） | 难度：中等
\begin{question}
    求 $\displaystyle\sum_{k=1}^{50}(3k+1)^{2}$ 的值。\paren[D]
    \begin{choices}
        \item $74000$
        \item $202775$
        \item $206550$
        \item $394025$
    \end{choices}
\end{question}

\begin{solution}
    求 $\displaystyle\sum_{k=1}^{50}(3k+1)^{2}$。 \\
    展开后逐项利用 $\sum k^{2}$、$\sum k$ 及 $\sum 1$ 的公式求和。
    \begin{align*}
        \sum_{k=1}^{50}(3k+1)^{2} &= \sum_{k=1}^{50}\left(9k^{2}+6k+1\right) \\
        &= 9\sum_{k=1}^{50}k^{2}+6\sum_{k=1}^{50}k+\sum_{k=1}^{50}1
    \end{align*}
    代入 $\displaystyle\sum_{k=1}^{50}k^{2}=\frac{50\times51\times101}{6}=42925$，$\displaystyle\sum_{k=1}^{50}k=\frac{50\times51}{2}=1275$：
    \begin{align*}
        \sum_{k=1}^{50}(3k+1)^{2} &= 9(42925)+6(1275)+50 \\
        &= 386325+7650+50 \\
        &= 394025
    \end{align*}
    故选 D。
\end{solution}
%endregion

%region 范围：2026年度金宝培元独中毕业考·高中组·高级数学试卷一 | 考点：反函数与不动点方程 | 难度：中等
\begin{question}
    已知函数 $g(x)=\dfrac{1}{x}+3$，$x\neq0$，求方程 $g^{-1}(x)=x$ 的解。\paren[D]
    \begin{choices}
        \item $1$
        \item $0$ 或 $1$
        \item $\dfrac{3\pm\sqrt{5}}{2}$
        \item $\dfrac{3\pm\sqrt{13}}{2}$
    \end{choices}
\end{question}

\begin{solution}
    求满足 $g^{-1}(x)=x$ 的 $x$。 \\
    互为反函数的图象关于直线 $y=x$ 对称，故 $g^{-1}(x)=x$ 的解等价于 $g(x)=x$ 的解。
    \begin{align*}
        \frac{1}{x}+3 &= x \\
        1+3x &= x^{2} \\
        x^{2}-3x-1 &= 0
    \end{align*}
    由求根公式：
    \[
        x = \frac{3\pm\sqrt{13}}{2}
    \]
    故选 D。
\end{solution}
%endregion

% %region 范围：2026年度金宝培元独中毕业考·高中组·高级数学试卷一 | 考点：逆矩阵、伴随矩阵与余子式 | 难度：困难
% \begin{question}
%     已知矩阵 $A=\begin{pmatrix} 1 & 2 & a \\ 2 & b & 3 \\ c & 0 & 8 \end{pmatrix}$ 的逆矩阵为 $A^{-1}=\begin{pmatrix} -40 & u & x \\ 13 & v & y \\ 5 & w & z \end{pmatrix}$，求 $b$ 的值。\paren[D]
%     \begin{choices}
%         \item $-5$
%         \item $-1$
%         \item $1$
%         \item $5$
%     \end{choices}
% \end{question}

% \begin{solution}
%     已知 $A^{-1}$ 的第一列，求 $b$。 \\
%     利用 $A^{-1}=\dfrac{1}{\det A}\operatorname{adj}A$，第一列元素对应第一行的三个余子式除以行列式。
%     \begin{align*}
%         C_{11} &= \begin{vmatrix} b & 3 \\ 0 & 8 \end{vmatrix} = 8b \\
%         C_{12} &= -\begin{vmatrix} 2 & 3 \\ c & 8 \end{vmatrix} = 3c-16 \\
%         C_{13} &= \begin{vmatrix} 2 & b \\ c & 0 \end{vmatrix} = -bc
%     \end{align*}
%     由 $A^{-1}$ 第一列条件：
%     \begin{align*}
%         \frac{C_{11}}{\det A} &= -40 \ \Rightarrow \ b=-5\det A \\
%         \frac{C_{12}}{\det A} &= 13 \ \Rightarrow \ c=\frac{13\det A+16}{3} \\
%         \frac{C_{13}}{\det A} &= 5 \ \Rightarrow \ bc=-5\det A
%     \end{align*}
%     代入 $b=-5\det A$ 至 $bc=-5\det A$，得 $c=1$；代回第二式求得：
%     \begin{align*}
%         13\det A+16 &= 3 \\
%         \det A &= -1
%     \end{align*}
%     故 $b=-5\det A=-5\times(-1)=5$。
%     故选 D。
% \end{solution}
% %endregion

%region 范围：2026年度金宝培元独中毕业考·高中组·高级数学试卷一 | 考点：对数换底公式 | 难度：中等
\begin{question}
    若 $\log_{7}a=\dfrac{1}{\log_{121}49}$，则 $a$ 满足以下哪个不等式？\paren[B]
    \begin{choices}
        \item $0<a<10$
        \item $10<a<20$
        \item $20<a<100$
        \item $100<a<1000$
    \end{choices}
\end{question}

\begin{solution}
    已知 $\log_{7}a=\dfrac{1}{\log_{121}49}$，求 $a$ 所在的范围。 \\
    先化简右式，再用换底公式求出 $a$ 的值。
    \begin{align*}
        \log_{121}49 &= \frac{\ln 7^{2}}{\ln 11^{2}} = \frac{\ln 7}{\ln 11} = \log_{11}7
    \end{align*}
    因此：
    \begin{align*}
        \log_{7}a &= \frac{1}{\log_{11}7} = \log_{7}11 \\
        a &= 11
    \end{align*}
    由于 $10<11<20$，故选 B。
\end{solution}
%endregion

%region 范围：2026年度金宝培元独中毕业考·高中组·高级数学试卷一 | 考点：一元二次方程判别式 | 难度：基础
\begin{question}
    已知方程式 $x^{2}-(2k+2)x+k^{2}+1=0$ 有两个相异的实根，求 $k$ 的取值范围。\paren[A]
    \begin{choices}
        \item $k>0$
        \item $k\geq0$
        \item $k<0$
        \item $k>-\dfrac{1}{2}$
    \end{choices}
\end{question}

\begin{solution}
    求使方程有两个相异实根的 $k$ 的取值范围。 \\
    两根相异等价于判别式大于零。
    \begin{align*}
        \Delta &= (2k+2)^{2}-4(k^{2}+1) \\
        &= 4k^{2}+8k+4-4k^{2}-4 \\
        &= 8k
    \end{align*}
    由 $\Delta>0$：
    \[
        8k>0 \ \Rightarrow \ k>0
    \]
    故选 A。
\end{solution}
%endregion

%region 范围：2026年度金宝培元独中毕业考·高中组·高级数学试卷一 | 考点：诱导公式与同角三角函数关系 | 难度：中等
\begin{question}
    若 $a>0$，$b>0$ 且 $\sin20^{\circ}=\dfrac{a}{b}$，则 $\tan160^{\circ}=?$\paren[D]
    \begin{choices}
        \item $\dfrac{a}{\sqrt{a^{2}+b^{2}}}$
        \item $-\dfrac{b}{\sqrt{a^{2}+b^{2}}}$
        \item $-\dfrac{b}{\sqrt{a^{2}-b^{2}}}$
        \item $-\dfrac{a}{\sqrt{b^{2}-a^{2}}}$
    \end{choices}
\end{question}

\begin{solution}
    已知 $\sin20^{\circ}=\dfrac{a}{b}$，求 $\tan160^{\circ}$。 \\
    利用诱导公式 $\tan160^{\circ}=-\tan20^{\circ}$，再由同角关系求出 $\cos20^{\circ}$。
    \begin{align*}
        \cos20^{\circ} &= \sqrt{1-\sin^{2}20^{\circ}} = \sqrt{1-\frac{a^{2}}{b^{2}}} = \frac{\sqrt{b^{2}-a^{2}}}{b}
    \end{align*}
    因此：
    \begin{align*}
        \tan20^{\circ} &= \frac{\sin20^{\circ}}{\cos20^{\circ}} = \frac{a}{\sqrt{b^{2}-a^{2}}} \\
        \tan160^{\circ} &= -\tan20^{\circ} = -\frac{a}{\sqrt{b^{2}-a^{2}}}
    \end{align*}
    故选 D。
\end{solution}
%endregion

%region 范围：2026年度金宝培元独中毕业考·高中组·高级数学试卷一 | 考点：余弦定理与三角形面积公式 | 难度：中等
\begin{question}
    在 $\triangle ABC$ 中，$\angle B$ 是钝角，$AB=6$，$BC=8$，$\triangle ABC$ 的面积为 $12\sqrt{3}$，求 $AC$ 的长。\paren[C]
    \begin{choices}
        \item $4.11$
        \item $7.21$
        \item $12.17$
        \item $13.53$
    \end{choices}
\end{question}

\begin{solution}
    已知 $AB=6$，$BC=8$，$\triangle ABC$ 面积为 $12\sqrt{3}$，$\angle B$ 为钝角，求 $AC$。 \\
    先由面积公式求 $\sin B$，再由钝角条件确定 $\cos B$，最后用余弦定理求 $AC$。
    \begin{align*}
        \frac{1}{2}\times AB\times BC\times\sin B &= 12\sqrt{3} \\
        24\sin B &= 12\sqrt{3} \\
        \sin B &= \frac{\sqrt{3}}{2}
    \end{align*}
    由 $\angle B$ 为钝角，取 $B=120^{\circ}$，故 $\cos B=-\dfrac{1}{2}$。
    \begin{align*}
        AC^{2} &= AB^{2}+BC^{2}-2\times AB\times BC\times\cos B \\
        &= 36+64-2(6)(8)\left(-\frac{1}{2}\right) \\
        &= 148 \\
        AC &= \sqrt{148} \approx 12.17
    \end{align*}
    故选 C。
\end{solution}
%endregion

%region 范围：2026年度金宝培元独中毕业考·高中组·高级数学试卷一 | 考点：三角方程与定义域限制 | 难度：中等
\begin{question}
    求方程式 $\tan x=2\sin x$ 在区间 $(0,2\pi)$ 内的解的个数。\paren[B]
    \begin{choices}
        \item $4$
        \item $3$
        \item $2$
        \item $1$
    \end{choices}
\end{question}

\begin{solution}
    求 $\tan x=2\sin x$ 在 $(0,2\pi)$ 内解的个数。 \\
    化为关于 $\sin x$ 的因式方程，注意 $\cos x\neq0$ 的定义域限制。
    \begin{align*}
        \frac{\sin x}{\cos x} &= 2\sin x \\
        \sin x\left(\frac{1}{\cos x}-2\right) &= 0
    \end{align*}
    情况一：$\sin x=0$，在 $(0,2\pi)$ 内得 $x=\pi$（且 $\cos\pi\neq0$，合法）。

    情况二：$\cos x=\dfrac{1}{2}$，在 $(0,2\pi)$ 内得 $x=\dfrac{\pi}{3}$ 或 $x=\dfrac{5\pi}{3}$。

    综合两种情况，共有 $3$ 个解。
    故选 B。
\end{solution}
%endregion

%region 范围：2026年度金宝培元独中毕业考·高中组·高级数学试卷一 | 考点：圆的标准式与直线相切条件 | 难度：中等
\begin{question}
    若直线 $3x+4y+k=0$ 与圆 $x^{2}+y^{2}+4x-6y-12=0$ 相切，求 $k$ 的值。\paren[C]
    \begin{choices}
        \item $-31$ 或 $-19$
        \item $31$ 或 $-19$
        \item $-31$ 或 $19$
        \item $31$ 或 $19$
    \end{choices}
\end{question}

\begin{solution}
    已知直线 $3x+4y+k=0$ 与已知圆相切，求 $k$。 \\
    先将圆化为标准式求出圆心与半径，再令圆心到直线的距离等于半径。
    \begin{align*}
        x^{2}+4x+y^{2}-6y &= 12 \\
        (x+2)^{2}+(y-3)^{2} &= 25
    \end{align*}
    圆心为 $(-2,3)$，半径为 $5$。由相切条件：
    \begin{align*}
        \frac{|3(-2)+4(3)+k|}{5} &= 5 \\
        |k+6| &= 25
    \end{align*}
    解得：
    \[
        k=19 \ \text{或} \ k=-31
    \]
    故选 C。
\end{solution}
%endregion

%region 范围：2026年度金宝培元独中毕业考·高中组·高级数学试卷一 | 考点：三角形面积公式（坐标法） | 难度：中等
\begin{question}
    已知 $A(k+1,6)$，$B(3,2)$ 及 $C(-k,0)$ 是 $\triangle ABC$ 的顶点。若 $\triangle ABC$ 的面积是 $10$ 平方单位，求 $k$ 的值。\paren[D]
    \begin{choices}[columns=2,linesep=0pt]
        \item $k=2$
        \item $k=6$
        \item $k=6$ 或 $k=-\dfrac{2}{3}$
        \item $k=2$ 或 $k=-18$
    \end{choices}
\end{question}

\begin{solution}
    已知三顶点坐标，面积为 $10$，求 $k$。 \\
    代入三角形面积的坐标公式，化简后取绝对值等于 $10$。
    \begin{align*}
        10 &= \frac{1}{2}\left|(k+1)(2-0)+3(0-6)+(-k)(6-2)\right| \\
        20 &= \left|2k+2-18-4k\right| \\
        20 &= \left|-2k-16\right| \\
        10 &= |k+8|
    \end{align*}
    解得：
    \[
        k=2 \ \text{或} \ k=-18
    \]
    故选 D。
\end{solution}
%endregion

% %region 范围：2026年度金宝培元独中毕业考·高中组·高级数学试卷一 | 考点：中位数与平均数 | 难度：基础
% \begin{question}
%     下列为一组顺序排列的数据：$3,3,4,7,x,12,14,18$。若数据的中位数与平均数相等，求 $x$ 的值。\paren[A]
%     \begin{choices}
%         \item $11$
%         \item $10$
%         \item $9$
%         \item $8$
%     \end{choices}
% \end{question}

% \begin{solution}
%     已知 $8$ 个按序排列的数据，中位数等于平均数，求 $x$。 \\
%     中位数为中间两数之平均，平均数为总和除以 $8$，两者相等列方程求解。
%     \begin{align*}
%         \text{中位数} &= \frac{7+x}{2}, \qquad \text{平均数} = \frac{61+x}{8}
%     \end{align*}
%     \begin{align*}
%         \frac{7+x}{2} &= \frac{61+x}{8} \\
%         4(7+x) &= 61+x \\
%         28+4x &= 61+x \\
%         3x &= 33 \\
%         x &= 11
%     \end{align*}
%     故选 A。
% \end{solution}
% %endregion

%region 范围：2026年度金宝培元独中毕业考·高中组·高级数学试卷一 | 考点：超几何分布（组合概率） | 难度：中等
\begin{question}
    一袋中盛有 $3$ 粒白球、$4$ 粒红球与 $5$ 粒黑球。若从袋中任意抽出三球，求三球中恰有二球是白球的概率。\paren[D]
    \begin{choices}
        \item $\dfrac{3}{64}$
        \item $\dfrac{9}{64}$
        \item $\dfrac{9}{220}$
        \item $\dfrac{27}{220}$
    \end{choices}
\end{question}

\begin{solution}
    从 $12$ 粒球中抽 $3$ 粒，求恰有 $2$ 粒白球的概率。 \\
    利用超几何分布，分子为从白球中选 $2$ 粒、其余球中选 $1$ 粒的组合数。
    \begin{align*}
        P &= \frac{\dbinom{3}{2}\dbinom{9}{1}}{\dbinom{12}{3}} \\
        &= \frac{3\times9}{220} \\
        &= \frac{27}{220}
    \end{align*}
    故选 D。
\end{solution}
%endregion

% %region 范围：2026年度金宝培元独中毕业考·高中组·高级数学试卷一 | 考点：含重复元素的排列（相邻元素不相邻问题） | 难度：困难
% \begin{question}
%     将 "coefficient" 一字的字母全取来排列，若两个 "f" 不能相邻，有多少种排法？\paren[D]
%     \begin{choices}
%         \item $36288000$
%         \item $4536000$
%         \item $4082400$
%         \item $2041200$
%     \end{choices}
% \end{question}

% \begin{solution}
%     求 "coefficient" 全部字母排列中两个 "f" 不相邻的排法数。 \\
%     先求全部排列数，再减去两个 "f" 相邻的排列数。
%     \begin{align*}
%         \text{总排列数} &= \frac{11!}{2!\,2!\,2!\,2!} = 2494800
%     \end{align*}
%     （字母 c, e, f, i 各重复 $2$ 次，共 $11$ 个字母。）
%     将两个 "f" 捆绑为一个整体，与其余 $9$ 个字母共 $10$ 个单位排列：
%     \begin{align*}
%         \text{相邻排列数} &= \frac{10!}{2!\,2!\,2!} = 453600
%     \end{align*}
%     （此时 c, e, i 各重复 $2$ 次。）
%     \begin{align*}
%         \text{不相邻排列数} &= 2494800-453600 = 2041200
%     \end{align*}
%     故选 D。
% \end{solution}
% %endregion

%region 范围：2026年度金宝培元独中毕业考·高中组·高级数学试卷一 | 考点：高阶导数的规律推导 | 难度：中等
\begin{question}
    已知 $f(x)=\dfrac{1}{x}$，求 $f^{(10)}(x)$。\paren[A]
    \begin{choices}
        \item $\dfrac{10!}{x^{11}}$
        \item $-\dfrac{10!}{x^{10}}$
        \item $\dfrac{10!}{x^{10}}$
        \item $\dfrac{11!}{x^{11}}$
    \end{choices}
\end{question}

\begin{solution}
    求 $f(x)=x^{-1}$ 的第十阶导数。 \\
    逐阶求导找出规律 $f^{(n)}(x)=(-1)^{n}n!\,x^{-(n+1)}$。
    \begin{align*}
        f'(x) &= -x^{-2} \\
        f''(x) &= 2x^{-3} \\
        f'''(x) &= -6x^{-4} \\
        f^{(n)}(x) &= (-1)^{n}n!\,x^{-(n+1)}
    \end{align*}
    代入 $n=10$：
    \[
        f^{(10)}(x) = (-1)^{10}\cdot10!\cdot x^{-11} = \frac{10!}{x^{11}}
    \]
    故选 A。
\end{solution}
%endregion

%region 范围：2026年度金宝培元独中毕业考·高中组·高级数学试卷一 | 考点：隐函数微分与切线斜率 | 难度：中等
\begin{question}
    求曲线 $x^{3}-xy+y^{2}=3$ 在点 $(1,2)$ 处的切线斜率。\paren[A]
    \begin{choices}
        \item $-\dfrac{1}{3}$
        \item $-\dfrac{5}{3}$
        \item $-\dfrac{13}{5}$
        \item $-1$
    \end{choices}
\end{question}

\begin{solution}
    求曲线在 $(1,2)$ 处的切线斜率。 \\
    对方程两边关于 $x$ 隐函数求导，再代入该点求 $\dfrac{dy}{dx}$。
    \begin{align*}
        3x^{2}-\left(y+x\frac{dy}{dx}\right)+2y\frac{dy}{dx} &= 0 \\
        (2y-x)\frac{dy}{dx} &= y-3x^{2}
    \end{align*}
    代入 $x=1,\ y=2$：
    \begin{align*}
        (4-1)\frac{dy}{dx} &= 2-3 \\
        \frac{dy}{dx} &= -\frac{1}{3}
    \end{align*}
    故选 A。
\end{solution}
%endregion

%region 范围：2026年度金宝培元独中毕业考·高中组·高级数学试卷一 | 考点：三角函数定积分（降幂公式） | 难度：中等
\begin{question}
    求 $\displaystyle\int_{0}^{\frac{\pi}{4}}\sin^{2}x\,dx$。\paren[A]
    \begin{choices}
        \item $\dfrac{\pi}{8}-\dfrac{1}{4}$
        \item $\dfrac{\pi}{8}$
        \item $\dfrac{\pi}{4}-\dfrac{1}{2}$
        \item $\dfrac{\pi}{4}$
    \end{choices}
\end{question}

\begin{solution}
    求 $\displaystyle\int_{0}^{\frac{\pi}{4}}\sin^{2}x\,dx$。 \\
    用降幂公式 $\sin^{2}x=\dfrac{1-\cos2x}{2}$ 化简后积分。
    \begin{align*}
        \int_{0}^{\frac{\pi}{4}}\sin^{2}x\,dx &= \int_{0}^{\frac{\pi}{4}}\frac{1-\cos2x}{2}\,dx \\
        &= \left[\frac{x}{2}-\frac{\sin2x}{4}\right]_{0}^{\frac{\pi}{4}}
    \end{align*}
    代入上下限：
    \begin{align*}
        &= \left(\frac{\pi}{8}-\frac{1}{4}\right)-(0-0) \\
        &= \frac{\pi}{8}-\frac{1}{4}
    \end{align*}
    故选 A。
\end{solution}
%endregion

% %region 范围：2026年度金宝培元独中毕业考·高中组·高级数学试卷一 | 考点：自反函数（involution）的积分 | 难度：困难
% \begin{question}
%     已知 $f(x)=\dfrac{x}{\sqrt{2x^{2}-1}}$，$x\geq\dfrac{1}{\sqrt2}$。求 $\displaystyle\int_{1}^{5}2f^{-1}(x)\,dx$。\paren[B]
%     \begin{choices}
%         \item $3$
%         \item $6$
%         \item $7$
%         \item $\dfrac{8}{7}$
%     \end{choices}
% \end{question}

% \begin{solution}
%     求 $\displaystyle\int_{1}^{5}2f^{-1}(x)\,dx$，其中 $f(x)=\dfrac{x}{\sqrt{2x^{2}-1}}$。 \\
%     先证 $f$ 为自反函数（$f^{-1}=f$），再直接积分。
%     设 $y=\dfrac{x}{\sqrt{2x^{2}-1}}$，解出 $x$：
%     \begin{align*}
%         y^{2}(2x^{2}-1) &= x^{2} \\
%         x^{2}(2y^{2}-1) &= y^{2} \\
%         x &= \frac{y}{\sqrt{2y^{2}-1}}
%     \end{align*}
%     可见反函数与原函数同式，即 $f^{-1}(x)=f(x)=\dfrac{x}{\sqrt{2x^{2}-1}}$。于是：
%     \begin{align*}
%         \int_{1}^{5}2f^{-1}(x)\,dx &= \int_{1}^{5}\frac{2x}{\sqrt{2x^{2}-1}}\,dx \\
%         &= \Big[\sqrt{2x^{2}-1}\Big]_{1}^{5} \\
%         &= \sqrt{49}-\sqrt{1} \\
%         &= 7-1=6
%     \end{align*}
%     故选 B。
% \end{solution}
% %endregion

\newpage
\vfill % 与上方内容保持距离

\noindent
\newlength{\sigwidth}
\settowidth{\sigwidth}{（\hspace{4em}师）}

\begin{tabular*}{\textwidth}{@{\extracolsep{\fill}} l l l}
    出题： & 审题： & 审核： \\[4em]
    \rule{\sigwidth}{0.5pt} &
    \rule{\sigwidth}{0.5pt} &
    \rule{\sigwidth}{0.5pt} \\
    （\hspace{0em} \setterteacher ） &
    （\hspace{0em} \reviewteacher ）  &
    （\hspace{0em} \auditteacher ） \\
\end{tabular*}

\clearpage
\thispagestyle{empty} % 取消这一页的页眉页脚（页码）
~
\clearpage
\end{document}
